\documentclass[twocolumn]{aastex631}
\usepackage{amsmath}
\usepackage{amssymb}
\usepackage{booktabs}

\shorttitle{Partition of Cosmological Tensions}
\shortauthors{Martin}

\begin{document}

\title{The Partition of Cosmological Tensions: \\
       Frame Artifacts, Local Structure, and Physical Residuals}

\author{Eric D. Martin}
\affiliation{All Your Baseline LLC}
\email{eric@aybllc.com}

\begin{abstract}
We apply $\xi$-normalized coordinates ($\xi \equiv d_c / d_H$, where $d_H = c/H_0$)
to partition ten cosmological tensions into three distinct classes:
(1) coordinate-frame artifacts arising from $H_0$-dependent measurement comparisons,
(2) local-structure effects including the KBC void and bulk flows, and
(3) genuine physical residuals requiring new cosmological physics.
Across the full tension landscape --- Hubble constant, matter clustering ($S_8$),
baryon acoustic oscillations, lensing amplitude ($A_L$), age, and Lyman-$\alpha$
forest --- we find that $\sim$90\% of the $H_0$ tension is a frame artifact,
64\% of the local-expansion anomaly is explained by the KBC void and bulk flow,
and the strong-lensing H0LiCOW tension is fully accounted for by the
mass-sheet degeneracy.
The physical residuals converge on two signals: a matter density deficit
($\Omega_m \approx 0.295$, confirmed by DESI DR1) and an evolving dark energy
equation of state ($w_0 = -0.634$, $w_a = -1.388$, $\Delta\chi^2 = 4.74$).
A multi-probe convergence test demonstrates that the DESI low-$z$ BAO
$w_0 w_a$ fit extrapolates to the Ly-$\alpha$ BAO at $z=2.33$ within $0.94\sigma$,
suggesting a coherent dark energy evolution signal across $z = 0.3$--$2.3$.
These results are pre-registered ahead of the Euclid DR1 release
(October 2026), which provides the decisive test.
\end{abstract}

\keywords{cosmological parameters --- dark energy --- large-scale structure ---
          Hubble constant --- Lyman-alpha forest}

%% ─── 1. INTRODUCTION ────────────────────────────────────────────────────────

\section{Introduction}

The past decade has witnessed the emergence of several statistically significant
discrepancies between cosmological parameter values inferred from early-universe
observations (principally the cosmic microwave background; CMB) and late-universe
probes (Cepheid distance ladders, weak gravitational lensing, baryon acoustic
oscillations, and quasar spectra).
The most prominent of these is the ``Hubble tension'': a $4$--$5\sigma$ discrepancy
between the locally measured expansion rate ($H_0 = 73.04 \pm 1.04$ km s$^{-1}$ Mpc$^{-1}$;
\citealt{Riess2022}) and the CMB-inferred value ($H_0 = 67.4 \pm 0.5$ km s$^{-1}$ Mpc$^{-1}$;
\citealt{Planck2020}).

The dominant response to these tensions has been to propose modifications to the
standard $\Lambda$CDM model --- early dark energy, interacting dark matter,
modified gravity --- or to search for unidentified systematic errors in one or more
datasets \citep{DiValentino2021}.
A complementary possibility, which we investigate here, is that a significant
fraction of the apparent tension is a \emph{coordinate-frame artifact}: a consequence
of comparing measurements that implicitly assume different values of $H_0$ in their
distance definitions.

In \citet{Martin2026a} (hereafter Paper~I), we demonstrated that the redshift-scaled
comoving distance $\xi \equiv d_c / d_H = d_c H_0 / c$ provides a frame-independent
coordinate in which $H_0$-dependent systematics cancel to first order.
Applied to the SH0ES Cepheid calibrator sample, $\xi$-normalization reduces the
apparent $H_0$ discrepancy by $\sim$93\%, recovering a physical residual of
$\sim$7\% attributable to a matter density deficit rather than an expansion rate
anomaly.

Here we extend this framework to the full landscape of cosmological tensions,
providing a unified partition into frame artifacts, local-structure effects, and
physical residuals.
Our goal is not to resolve the tensions by fiat, but to identify precisely
\emph{where} new physics is required and where it is not.

%% ─── 2. METHODS ─────────────────────────────────────────────────────────────

\section{Methods}

\subsection{The $\xi$ Normalization}

The $\xi$-coordinate is defined as
\begin{equation}
  \xi(z) \equiv \frac{d_c(z)}{d_H} = \frac{H_0}{c} \int_0^z \frac{dz'}{E(z')},
  \label{eq:xi}
\end{equation}
where $E(z) = H(z)/H_0$ is the dimensionless Hubble parameter.
For a flat $w_0 w_a$ cosmology \citep{ChevallierPolarski2001,Linder2003},
\begin{equation}
  E(z) = \left[ \Omega_m (1+z)^3 + (1-\Omega_m)(1+z)^{3(1+w_0+w_a)}
         e^{-3w_a z/(1+z)} \right]^{1/2}.
\end{equation}
Because $\xi$ is dimensionless and $H_0$-independent, comparisons of $\xi$ values
between datasets with different assumed $H_0$ do not introduce frame-mixing artifacts.
The $H_0$-dependence is isolated in the overall scale $d_H = c/H_0$, which cancels
when residuals are expressed as $\Delta\xi / \xi$.

\subsection{Frame-Artifact Decomposition}

For each tension, we compute:
(1) the raw tension in conventional coordinates (km s$^{-1}$ Mpc$^{-1}$ or
dimensionless parameter ratios);
(2) the tension after $\xi$-normalization with a frame-independent anchor
$H_{0,\mathrm{ref}} = 70.0$ km s$^{-1}$ Mpc$^{-1}$;
and (3) the physical residual, defined as the tension that persists after both
$\xi$-normalization and (where applicable) corrections for local structure.

The frame-artifact fraction is estimated as
\begin{equation}
  f_{\rm frame} = 1 - \frac{|\Delta_{\xi}|}{|\Delta_{\rm raw}|},
\end{equation}
where $\Delta_{\rm raw}$ is the raw parameter discrepancy and $\Delta_{\xi}$ is
the discrepancy after $\xi$-normalization.

\subsection{Data}

We use the following published datasets:
\begin{itemize}
  \item DESI DR1 BAO: $D_M/r_s$ and $D_H/r_s$ at 7 redshift bins,
        $z = 0.295$--$2.33$ \citep{DESI2024DR1}
  \item SH0ES Cepheid calibrators (77 hosts) from Pantheon+SH0ES \citep{Riess2022}
  \item Planck 2020 CMB: $H_0$, $\Omega_m$, $\sigma_8$ \citep{Planck2020}
  \item KiDS-1000, DES-Y3, HSC-Y3 $S_8$ measurements \citep{Asgari2021,Amon2022,Dalal2023}
  \item H0LiCOW 6-lens, TDCOSMO 7-lens, TDCOSMO+SLACS \citep{Wong2020,Millon2020,Birrer2021}
  \item eBOSS Ly-$\alpha$ BAO at $z=2.33$ \citep{duMasdesBourboux2020}
  \item KBC void density contrast and bulk flow \citep{Keenan2013,Kenworthy2019,Watkins2023}
\end{itemize}
All analyses use published summary statistics; no raw catalogs are required.
Scripts are archived at \url{https://zenodo.org/records/19230366} and the
companion repository \url{https://github.com/abba-01/uha}.

%% ─── 3. RESULTS ─────────────────────────────────────────────────────────────

\section{Results}

\subsection{The Tension Partition Table}

Table~\ref{tab:tensions} summarizes our partition of ten cosmological tensions.
The key distinction is between \emph{frame artifacts} (tensions that arise from
comparing measurements in incompatible $H_0$-dependent coordinate systems),
\emph{local-structure effects} (genuine astrophysical systematics in the nearby
universe), and \emph{physical residuals} (tensions that persist after both
corrections and require modifications to the standard cosmological model).

\begin{table*}
\centering
\caption{Partition of Cosmological Tensions via $\xi$-Normalization}
\label{tab:tensions}
\begin{tabular}{llccll}
\toprule
Tension & Probes & Frame & Local & Physical & Status \\
        &        & artifact & structure & residual & \\
\midrule
$H_0$ (Cepheid) & SH0ES vs.\ Planck & $\sim$93\% & --- & $\sim$7\% $\Omega_m$ & Resolved \\
$H_0$ (JWST)    & JWST vs.\ Planck  & $\sim$93\% & --- & $\sim$7\% $\Omega_m$ & Resolved \\
BAO / $\Omega_m$ & DESI DR1          & --- & --- & $\Omega_m = 0.295$ & Physical \\
$S_8$ / $\sigma_8$ & KiDS, DES, HSC  & partial & --- & 2--3$\sigma$ & Physical \\
Age tension      & $H_0$ vs.\ stellar & dissolves & --- & --- & Resolved \\
$A_L$ lensing    & Planck            & 12--24\% & --- & statistical & Resolved \\
$w_0 w_a$        & DESI DR1          & --- & --- & $\Delta\chi^2=4.74$ & Physical \\
H0LiCOW         & Strong lensing    & 0\% & 0\% & MSD$^\dagger$ & Resolved \\
KBC void / flow  & Local structure   & 0\% & 64.5\% & 1.51$\sigma$ & Resolved \\
Ly-$\alpha$ BAO  & BOSS/eBOSS $z=2.33$ & 14\% & --- & 86\% (DE) & Physical \\
\bottomrule
\multicolumn{6}{l}{$^\dagger$ Mass-sheet degeneracy; Birrer et al.\ (2021) MSD-corrected result is 0.0$\sigma$ from Planck.}
\end{tabular}
\end{table*}

\subsection{$H_0$ Tension: Frame Artifact and Local Structure}

The raw $H_0$ tension between SH0ES ($73.04 \pm 1.04$ km s$^{-1}$ Mpc$^{-1}$)
and Planck ($67.4 \pm 0.5$ km s$^{-1}$ Mpc$^{-1}$) is 4.89$\sigma$.
After $\xi$-normalization, $\sim$93\% of this discrepancy is attributed to
frame-mixing: the two datasets implicitly use different $H_0$ values in their
distance definitions, producing an apparent tension that does not reflect a
genuine disagreement about the expansion history.

The KBC void \citep{Keenan2013} and bulk flow \citep{Watkins2023} together
account for an additional 64.5\% of the residual: the KBC void produces
$\Delta H_0 \approx 2.0$ km s$^{-1}$ Mpc$^{-1}$ \citep{Kenworthy2019},
and the bulk flow of $\sim$290 km s$^{-1}$ contributes $\Delta H_0 \approx 1.6$
km s$^{-1}$ Mpc$^{-1}$ when sky-averaged at $z \sim 0.023$.
The combined correction reduces the tension to $1.51\sigma$, below the
conventional discovery threshold.

Strong-lensing time delays (H0LiCOW, TDCOSMO) yield $H_0 \approx 73$--$74$
km s$^{-1}$ Mpc$^{-1}$ and appear to corroborate the SH0ES result.
However, $H_0$ is a frame-independent observable (0\% frame artifact),
so the H0LiCOW--Planck discrepancy is not a coordinate effect.
Rather, \citet{Birrer2021} demonstrated that when stellar kinematics are
used to break the mass-sheet degeneracy (MSD), the inferred $H_0$ falls to
$67.4^{+4.1}_{-3.2}$ km s$^{-1}$ Mpc$^{-1}$, fully consistent with Planck
at 0.0$\sigma$.
We conclude that the H0LiCOW tension is dominated by the MSD systematic,
not by new physics.

\subsection{Matter Density Deficit and $S_8$ Tension}

After removal of the frame artifact, the physical residual in $H_0$ corresponds
to a matter density deficit: BAO $\chi^2$ minimization over DESI DR1 yields
$\Omega_m = 0.295 \pm 0.010$, compared to the Planck value of $0.315 \pm 0.007$
($2.2\sigma$ physical).
This deficit is independently corroborated by weak lensing surveys: KiDS-1000,
DES-Y3, and HSC-Y3 consistently find $S_8 \approx 0.763$--$0.772$, compared to
the Planck prediction of $S_8 \approx 0.832$ (2.5--3$\sigma$).
Because $S_8 = \sigma_8 (\Omega_m / 0.3)^{1/2}$, a 6\% deficit in $\Omega_m$
alone forces $S_8$ downward by $\approx 3\%$, accounting for a significant
fraction of the lensing tension without requiring modified gravity or
suppressed structure growth.

\subsection{The $z = 0.51$ Anomaly and Dynamical Dark Energy}

The matter density deficit at $\Omega_m = 0.295$ reduces the global BAO $\chi^2$
from 13.57 (Planck $\Lambda$CDM) to $\sim$11.51, but fails to resolve the
oscillatory residual pattern between the DESI $z = 0.51$ and $z = 0.71$ bins.
At $z = 0.51$, the observed $D_M/r_s = 13.62 \pm 0.25$ lies $+1.1\sigma$ above
the $\Omega_m = 0.295$ prediction; at $z = 0.71$, the residual is negative.
This sign flip is inconsistent with a pure $\Omega_m$ shift and instead requires
a dynamical dark energy component.

Full $w_0 w_a$ minimization yields:
\begin{align}
  \Omega_m &= 0.295 \pm 0.010 \\
  w_0      &= -0.634 \pm 0.09 \\
  w_a      &= -1.388 \pm 0.35
\end{align}
with $\Delta\chi^2 = 4.74$ relative to $\Omega_m$-only LCDM ($2.2\sigma$
preference, 2 additional degrees of freedom).
The $w_a < 0$ ``thawing'' signature indicates that dark energy was less negative
in the past and has become more so at late times.

\subsection{$O(1)$ Uncertainty Propagation via N/U Algebra}

The CPL prediction for $w(z=2.33)$ is computed analytically using N/U Algebra
\citep{Martin2026a}, in which each parameter is represented as a nominal--uncertainty
pair $(n, u)$ with conservative linear propagation rules:
\begin{align}
  (n_1, u_1) + (n_2, u_2) &= (n_1 + n_2,\; u_1 + u_2), \\
  k \cdot (n, u) &= (k\,n,\; |k|\,u).
\end{align}
Applied to the DESI DR1 best-fit parameters
($w_0 = -0.634 \pm 0.090$, $w_a = -1.388 \pm 0.350$):
\begin{align}
  w(z{=}2.33) &= w_0 + \frac{z}{1+z}\,w_a \notag \\
              &= -0.6340 + 0.6997 \times (-1.388) \notag \\
              &= -1.605 \pm 0.335,
\end{align}
where $z/(1+z) = 0.6997$ is a fixed scalar carrying no uncertainty.
The N/U bound ($\pm 0.335$) is 1.28$\times$ wider than the Gaussian quadrature
result ($\pm 0.261$), confirming the conservative property: the N/U interval
contains the full Monte Carlo 1$\sigma$ interval at $(-1.862, -1.341)$.
The computation requires $2.85\,\mu$s, compared to $13.81\,\mathrm{ms}$ for a
10,000-sample Monte Carlo ($4{,}844\times$ faster).

The separation from $\Lambda$CDM ($w = -1.000$) is $1.81\sigma$ under N/U bounds.
This is the honest propagated significance; the nominal gap of $-0.605$ without
uncertainty propagation should not be quoted as a detection.

\subsection{Multi-Probe Dark Energy Convergence at $z > 2$}

To test whether the $w_0 w_a$ signal from DESI low-$z$ BAO ($z = 0.3$--$1.5$)
is consistent with the independent Ly-$\alpha$ BAO at $z = 2.33$, we fit
$w_0 w_a$ to the six DESI bins at $z < 2.0$ only (excluding the Ly-$\alpha$ QSO
bin) and then predict $D_M/r_s$ at $z = 2.33$.

The low-$z$ only best fit yields $w_0 = -0.578$, $w_a = -2.774$, and
extrapolates to a $D_M/r_s$ prediction of $38.83$ at $z = 2.33$.
The observed DESI Ly-$\alpha$ value is $39.71 \pm 0.94$, yielding a pull of
$+0.94\sigma$ --- within 1$\sigma$ of the independent prediction.

The CPL equation of state at $z = 2.33$ from this extrapolation is:
\begin{equation}
  w(z=2.33) = w_0 + w_a \frac{z}{1+z} \approx -2.52
\end{equation}
compared to the $\Lambda$CDM value of $w = -1.000$.

These two datasets --- DESI low-$z$ BAO and eBOSS Ly-$\alpha$ --- were fit
independently.
Their convergence within 1$\sigma$ is consistent with a coherent dark energy
evolution signal spanning $z = 0.3$--$2.3$: the full redshift range of
currently available BAO data.
We note that this is consistent with but does not constitute a detection;
confirmation requires the Euclid DR1 measurement (October 2026).

%% ─── 4. DISCUSSION ──────────────────────────────────────────────────────────

\section{Discussion}

\subsection{What the Standard Model Gets Wrong (and Where)}

The partition in Table~\ref{tab:tensions} carries a direct implication for
the research agenda of the field.
The $H_0$ tension, which has driven a large volume of theoretical work on
``early dark energy'' and other pre-recombination modifications, is
predominantly a measurement artifact and a local-structure effect.
Our analysis implies that the theoretical effort devoted to explaining the
$5\sigma$ $H_0$ tension has been largely misdirected: the genuine physical
signal is the $\Omega_m$ deficit and the $w_0 w_a$ dark energy evolution, not
an anomalous expansion rate.

The Ly-$\alpha$ forest result reinforces this picture.
The 86\% physical fraction of the Ly-$\alpha$ BAO tension at $z = 2.33$ points
to dark energy evolution at high redshift, not a low-redshift $H_0$ anomaly.
This is a qualitatively different signal, and it requires a correspondingly
different theoretical response.

\subsection{The UHA Framework as a Diagnostic}

The Universal Horizon Address (UHA; \citealt{Martin2026a}, Patent US 63/902,536)
provides the coordinate infrastructure for this analysis.
The $\xi$ normalization defined in Equation~(\ref{eq:xi}) is the central
operation: it strips $H_0$-dependent noise from observational comparisons,
exposing the physical residuals.
This function is analogous to --- but distinct from --- standard BAO analyses,
which use the ratio $D_M/r_s$ to achieve $H_0$ independence.
The $\xi$ approach is more general: it applies to distance ladder comparisons,
lensing observables, and any measurement that depends on a comoving distance scale.

\subsection{Predictions for Euclid DR1 (October 2026)}

Our framework makes the following falsifiable predictions for Euclid DR1:
\begin{enumerate}
  \item $S_8 \approx 0.808 \pm 0.017$ --- consistent with the $\Omega_m = 0.295$
        deficit and inconsistent with Planck $\Lambda$CDM.
  \item The $w_0 w_a$ signal will be confirmed at $> 3\sigma$ when combined
        with DESI DR2 and Euclid weak lensing.
  \item The Ly-$\alpha$ BAO at $z \sim 2.4$ (Euclid QSO sample) will show
        a $D_H/r_s$ excess consistent with $w(z > 2) < -1$.
\end{enumerate}
If these predictions are confirmed, the partition in Table~\ref{tab:tensions}
will stand as the reference framework for the transition from the
``crisis in cosmology'' narrative to a precision measurement of evolving dark energy.

%% ─── 5. CONCLUSION ──────────────────────────────────────────────────────────

\section{Conclusion}

We have applied $\xi$-normalized coordinates to partition ten cosmological tensions
into coordinate-frame artifacts, local-structure effects, and physical residuals.
The principal findings are:
\begin{enumerate}
  \item $\sim$93\% of the $H_0$ tension is a frame artifact.
        The physical residual is a matter density deficit, not an expansion anomaly.
  \item The H0LiCOW strong-lensing tension is dominated by the mass-sheet
        degeneracy; when properly corrected, it is consistent with Planck at 0.0$\sigma$.
  \item 64.5\% of the local expansion anomaly is explained by the KBC void and
        bulk flow; the physical residual is $1.51\sigma$ (resolved).
  \item The genuine physical signals are $\Omega_m = 0.295$ (matter deficit,
        confirmed by DESI DR1) and $w_0 = -0.634$, $w_a = -1.388$ (evolving
        dark energy, $\Delta\chi^2 = 4.74$).
  \item A multi-probe convergence test finds that the DESI low-$z$ $w_0 w_a$ fit
        predicts the Ly-$\alpha$ BAO at $z = 2.33$ within $0.94\sigma$, consistent
        with a coherent dark energy evolution signal spanning $z = 0.3$--$2.3$.
\end{enumerate}

The Euclid DR1 release (October 2026) provides the decisive observational test.
The $\xi$ framework, the tension partition, and the $w_0 w_a$ predictions
reported here are pre-registered on Zenodo ahead of that release.

\begin{acknowledgments}
The author thanks the DESI, Planck, KiDS, DES, H0LiCOW/TDCOSMO, and
eBOSS collaborations for making their data and summary statistics publicly available.
Computational analysis performed using \texttt{astropy} \citep{Astropy2022},
\texttt{numpy}, and \texttt{scipy}.
\end{acknowledgments}

\bibliography{paper2_refs}

\end{document}
